\begin{description}
\item[2.1:] Figure 2 shows a part of the sky which contains \textbf{Cephei
  constellation}, for 22 October 2009 at 22:00 local time. Five bright stars
  in Cephei constellation are identified by numbers (1, 2, \dots, 5) and
  common names. Estimate the angular distances (in units of degrees) between
  two pairs of stars shown in table 2-1 and complete this table with your
  answers. \textbf{(40 Points)}

  \medskip
  \noindent Table 2-1: Angular Distance

  \begin{center}
  \begin{tabular}{lc}
  \hline
  Pairs of stars & Angular Distance (degrees) \\
  \hline
  1 (Errai) and 2 (Alfirk) & \\
  1 (Errai) and 3 (Alderamin) & \\
  \hline
  \end{tabular}
  \end{center}

\item[2.2:] Use table 2-2 and figure 2, then Estimate the ``apparent visual
  magnitude'' of stars 2 (Alfirak) and 3 (Alderamin) and complete Table 2-3.
  \textbf{(40 Points)}

  \medskip
  \noindent Table 2-2

  \begin{center}
  \begin{tabular}{cc}
  \hline
  Star Name & Apparent Visual Magnitude \\
  \hline
  Polaris & 1.95 \\
  Altais & 3.05 \\
  Segin & 3.34 \\
  \hline
  \multicolumn{2}{c}{All of these stars, are marked in the figure 2} \\
  \hline
  \end{tabular}
  \end{center}

  \medskip
  \noindent Table 2-3: Magnitude Estimation

  \begin{center}
  \begin{tabular}{ccc}
  \hline
  Star Number & Star Name & Apparent Visual Magnitude \\
  \hline
  2 & Alfirk & \\
  3 & Alderamin & \\
  \hline
  \end{tabular}
  \end{center}
\end{description}

\ioaafig{2009_OBS-SKY_2-f1.pdf}
% ($IOAA_PAPERS_DIR/observation/2009_OBS.pdf, page 7): sky chart around
% Cepheus with stars 1 (Errai), 2 (Alfirk), 3 (Alderamin), 5 (Iota Cephei),
% plus Polaris, Altais and Segin marked.
